\bookchapter{Глава 12: Аудиовизуальные развлечения и чтение}

\chapterintro{Смартфон уже практически заменил MP3-плеер, пульт от телевизора, Kindle и портативную игровую консоль. Но если честно, большинство людей используют лишь самые базовые функции этих приложений, и многие полезные возможности остаются незамеченными. В этой главе мы поговорим о том, как сделать мобильные развлечения удобнее и комфортнее.}

\booksection{Прослушивание музыки: какое приложение выбрать и как настроить комфортное звучание}

Сейчас функции музыкальных приложений во многом схожи, но у каждого есть свои акценты — правильный выбор сэкономит вам немало нервов:

\begin{itemize}
  \item \textbf{NetEase Cloud Music}: его сильная сторона — сообщество плейлистов и комментарии. Если вы хотите послушать «песни для дождливого дня» или «инструментальную музыку для отдыха после учёбы», просто введите запрос в поиск — и найдёте готовые плейлисты, собранные другими пользователями. Комментарии иногда бывают интереснее самих песен. Недостаток: не все песни имеют полные авторские права, некоторые треки со временем становятся недоступны.
  \item \textbf{QQ Music}: самая полная библиотека авторских прав — песни таких популярных исполнителей, как Джей Чоу и Mayday, есть практически только в QQ Music. Богатый выбор качества звука, подписчики могут включить lossless-качество. Если для вас важно «найти любую песню, какую захотите», QQ Music — самый надёжный выбор.
  \item \textbf{Apple Music}: лучший выбор для пользователей Apple. Чистый интерфейс, без рекламы, отличное качество звука, поддержка Lossless и Dolby Atmos. Однако его рекомендательный алгоритм уступает китайским приложениям в понимании предпочтений китайскоязычных пользователей, а культура плейлистов не так развита, как в NetEase Cloud. Подходит тем, кто ценит лаконичный опыт и не слишком заботится о социальных функциях.
  \item \textbf{Spotify}: самое большое количество пользователей в мире, признанно лучший рекомендательный алгоритм. Точность плейлистов «Discover Weekly» и «Daily Mix» поразительно высока — чем дольше пользуетесь, тем лучше Spotify вас понимает. Недостатки: много рекламы в бесплатной версии, неполное покрытие авторских прав на некоторые китайские песни. Требуется специальное сетевое подключение.
\end{itemize}

\textbf{Несколько советов по прослушиванию}: если вы используете Bluetooth-наушники, зайдите в настройки музыкального приложения и поищите «звуковые эффекты» или «эквалайзер» — небольшая настройка низких и высоких частот заметно улучшит звучание. Кроме того, потоковая музыка расходует много трафика — не забудьте включить в настройках «воспроизводить lossless только по Wi-Fi» или заранее загрузите любимые плейлисты для офлайн-прослушивания.

\booksection{Подкасты: «компаньон для ушей» в дороге}

Подкасты — это то, что кажется неважным тем, кто никогда их не слушал, но стоит начать — и обратного пути нет. По сути, это «Netflix-версия радиопередач»: кто-то говорит о технологиях, кто-то рассказывает историю, кто-то читает рассказы — есть контент на любой вкус.

\begin{itemize}
  \item \textbf{Приложение Xiaoyuzhou}: лучшее приложение для подкастов в Китае — чистый интерфейс, без лишних функций. Его «ежедневные рекомендации» и «выбор редакции» помогают открывать новые программы, а раздел комментариев очень активен. Лично я чаще всего слушаю «Hu Zuo Hu You» (культура и история) и «Bu He Shi Yi» (социальные темы) — каждый выпуск будто подслушиваешь разговор нескольких умных людей.
  \item \textbf{Apple Podcasts}: встроено в iPhone, не требует отдельной загрузки. Лучшая интеграция с CarPlay — очень удобно за рулём.
  \item \textbf{Советы по прослушиванию}: большинство приложений для подкастов поддерживают воспроизведение на скорости 1,2x или 1,5x — привыкнув, вы сэкономите массу времени. Я обычно слушаю на 1,2x на китайском и на 1,5x на английском. Также обязательно включите «автоматически загружать новые выпуски» и «автоматически удалять после прослушивания» — так телефон не будет забит прослушанными аудиофайлами.
\end{itemize}

\booksection{Просмотр видео: правильный подход от плеера до стриминга}

Многие не знают, что на телефоне можно воспроизводить гораздо больше форматов, чем предлагают стриминговые приложения. Хороший локальный плеер решает множество проблем.

\begin{itemize}
  \item \textbf{VLC Media Player}: бесплатный и с открытым исходным кодом, воспроизводит практически любые форматы — MKV, AVI, FLAC, даже образы дисков ISO. Если вы скачали видео, которое работает на компьютере, но не на телефоне — скорее всего, дело в формате, и установка VLC решит проблему. Он также поддерживает внешние субтитры: положите файл .srt в одну папку с видео, и VLC загрузит их автоматически.
  \item \textbf{nPlayer}: если хотите более изысканный опыт, этот платный плеер стоит своих денег. Красивый интерфейс, удобное управление жестами: свайп вверх-вниз по левой стороне регулирует яркость, по правой — громкость, свайп влево-вправо — перемотка. Поддерживает DTS и Dolby Audio — при просмотре блокбастеров создаёт гораздо более захватывающий эффект.
  \item \textbf{Bilibili (B-станция)}: уже не просто сайт для аниме — сейчас здесь невероятное количество образовательных ресурсов. Подготовка к экзаменам, программирование, фотография, кулинария — есть уроки на любую тему. Культура «даньму» (летающих комментариев) — душа Bilibili: когда смотришь видео с включёнными комментариями, создаётся ощущение, что целая группа людей смотрит вместе с тобой. Но при просмотре учебных видео советую их отключать — иначе легко отвлечься.
  \item \textbf{Трансляция на телевизор}: как бы ни был велик экран телефона, смотреть на телевизоре всё равно комфортнее. В одной сети Wi-Fi нажмите кнопку «TV» в правом верхнем углу видеоприложения — и картинка передастся на телевизор. Большинство смарт-телевизоров и проекторов это поддерживают, и опыт просмотра сериалов и фильмов становится намного лучше. Пользователи iPhone также могут использовать AirPlay для прямого зеркалирования всего экрана.
\end{itemize}

\booksection{Чтение электронных книг: на телефоне тоже можно читать с удовольствием}

Некоторые говорят, что читать с телефона вредно для глаз и легко отвлекаться — в этом есть доля правды. Но у мобильного чтения есть и очевидные преимущества: доступно всегда и везде, максимально использует каждую свободную минуту. Я прочитал немало книг с телефона — в метро, в очереди, перед сном.

\begin{itemize}
  \item \textbf{WeRead (WeChat Reading)}: сейчас лучшее приложение для чтения в Китае. Обширная библиотека — есть и изданные книги, и веб-романы. Модель «читательское время в обмен на книжные монеты» позволяет много читающим получать книги бесплатно, что особенно удобно для студентов. Функция «подчёркивания» показывает, какие отрывки отметили другие читатели — иногда видишь выделенную кем-то хорошую фразу, и это действительно вызывает отклик. Однако рейтинг читателей создаёт ненужную конкуренцию — легко начать «читать ради наращивания времени», поэтому советую его отключить.
  \item \textbf{Kindle App}: если у вас есть устройство Kindle, мобильное приложение Kindle синхронизирует прогресс чтения. Качество вёрстки лицензионных электронных книг Kindle намного лучше пиратских PDF, но китайский книжный магазин Kindle уже закрыт — нужно переключиться на зарубежный регион.
  \item \textbf{Apple Books (Книги)}: встроено в iPhone, подходит для чтения EPUB и PDF. Минималистичный интерфейс, приятная анимация перелистывания. Если вы скачали файл электронной книги из интернета, просто «поделиться в Книги» — и она импортируется.
  \item \textbf{Рекомендации по настройкам чтения}: какое бы приложение вы ни использовали, обязательно установите тёплый жёлтый фон (режим защиты глаз), увеличьте шрифт и межстрочный интервал. Для вечернего чтения включите «тёмный режим» — белый текст на чёрном фоне, глазам гораздо комфортнее. Ни в коем случае не читайте долго белым по чёрному мелким шрифтом при слабом освещении — усталость глаз накапливается.
\end{itemize}

\booksection{Мобильные игры: играть с удовольствием и без вреда для телефона}

Моё отношение к мобильным играм такое: играть можно, но не позволяйте игре контролировать вас. Разумно используйте игровые инструменты — играйте комфортнее и не в ущерб делам.

\begin{itemize}
  \item \textbf{Игровой режим / Игровое пространство}: эта функция есть у большинства современных телефонов. При включении телефон автоматически очищает фоновые процессы, блокирует уведомления и предотвращает случайные касания — представьте, что вы в разгаре командного боя и вдруг выскакивает входящий звонок; игровой режим защитит вас от таких ситуаций. Некоторые телефоны также поддерживают режим производительности / E-sports режим, заставляющий процессор работать на полную мощность, но ценой более быстрого нагрева и расхода батареи.
  \item \textbf{Как настроить графику и частоту кадров}: не выкручивайте всё на максимум бездумно. Сначала обеспечьте частоту кадров (плавность), потом гонитесь за качеством картинки. 60 fps гораздо комфортнее 30 fps, но 120 fps создаёт серьёзную нагрузку на телефон — устройствам среднего класса не нужно к этому стремиться. Если вы играете в соревновательные игры (например, Honor of Kings, PUBG), приоритет частоты кадров намного выше приоритета графики — плавное управление в сто раз важнее красивой травы.
  \item \textbf{Что делать, если телефон перегревается во время игры}: это нормально, но длительный перегрев сокращает срок службы батареи. Несколько способов охлаждения: снимите чехол для лучшего рассеивания тепла; уменьшите яркость экрана примерно до 70\% (яркость — мощный источник нагрева); не играйте во время зарядки (нагрев суммируется, и батарея умирает быстрее). Если телефон раскалился так, что его трудно держать в руке, немедленно выйдите из игры и дайте ему отдохнуть.
  \item \textbf{Инструменты для защиты от игровой зависимости}: «Экранное время» в iOS и «Цифровое благополучие» в Android позволяют установить дневной лимит времени для игровых приложений. Я установил себе максимум час в будни и два часа в выходные. Когда время истекает, иконка приложения становится серой и не открывается. Это ограничение действительно работает, а не просто для галочки.
  \item \textbf{Облачный гейминг}: если ваш телефон недостаточно мощный, но вы хотите играть в крупные проекты, попробуйте облачные игровые платформы (например, Tencent START Cloud Gaming, NetEase Cloud Gaming). Игра запускается на удалённом сервере, а изображение передаётся на ваш телефон в виде видеопотока. Требуется хорошее интернет-соединение, но к производительности телефона почти нет требований. Недостаток — неизбежная задержка; подходит для стратегий и RPG, но не для шутеров от первого лица.
\end{itemize}

\begin{tcbpractice}{Практическое задание: создайте свою систему мобильных развлечений}
\begin{enumerate}
  \item Загрузите 3 любимых альбома или плейлиста в музыкальном приложении и убедитесь, что они доступны офлайн.
  \item Подпишитесь на 3 интересных канала в Xiaoyuzhou или Apple Podcasts, загрузите один выпуск и послушайте.
  \item Установите VLC Media Player и попробуйте воспроизвести видеофайл, скопированный с компьютера.
  \item Найдите в WeRead книгу, которую хотите прочитать, установите тёплый фон и режим защиты глаз, настройте шрифт подходящего размера.
  \item Если вы играете в игры, откройте игровой режим и посмотрите, какие настройки можно оптимизировать (беззвучный режим, приоритет частоты кадров и т.д.).
\end{enumerate}

\end{tcbpractice}
